\documentclass[11pt]{article}

\usepackage[a4paper,margin=29mm]{geometry}
\usepackage[T1]{fontenc}
\usepackage{lmodern}
\usepackage{amsmath,amssymb,amsthm,mathtools}
\usepackage{booktabs}
\usepackage{enumitem}
\usepackage{xcolor}
\usepackage[hidelinks]{hyperref}
\usepackage{microtype}

\newtheorem{theorem}{Theorem}[section]
\newtheorem{lemma}[theorem]{Lemma}
\newtheorem{proposition}[theorem]{Proposition}
\newtheorem{corollary}[theorem]{Corollary}
\theoremstyle{definition}
\newtheorem{definition}[theorem]{Definition}
\newtheorem{remark}[theorem]{Remark}

\newcommand{\N}{\mathbb N}
\newcommand{\Q}{\mathbb Q}
\newcommand{\leanname}[1]{\nolinkurl{#1}}
\newcommand{\HSC}{\mathrm{HSC}}
\hypersetup{
  pdftitle={The Herzog-Schönheim conjecture for coset partitions with at most seventeen cells},
  pdfauthor={Rio Itabe}
}

\title{The Herzog--Sch\"onheim conjecture for coset partitions\\
with at most seventeen cells\\[3mm]
\large An unreviewed computer-assisted candidate}
\author{Rio Itabe\\\small Independent researcher}
\date{Version 0.4.2-review-candidate, 17 August 2026}

\begin{document}
\maketitle

\begin{center}
\small This manuscript has not been peer reviewed. It presents a
computer-assisted proof candidate for a bounded partial result only. The
unrestricted Herzog--Sch\"onheim conjecture remains open, and no claim of
priority is made.
\end{center}

\begin{abstract}
The Herzog--Sch\"onheim conjecture asserts that a nontrivial finite partition
of a group into left cosets of finite-index subgroups has two equal subgroup
indices. We present an unreviewed computer-assisted proof candidate that the
conjecture holds for every partition with at most seventeen cosets.
Equivalently, every counterexample
requires at least eighteen cells. After index-two descent and four published
Margolis--Schnabel obstructions, exact arithmetic enumeration leaves five
profiles at length seventeen, all beginning with index four. An index-four
cell forces a finite assignment to three non-anchor left-coset boxes. If a
cell of ambient index $n$ meets $d$ boxes and has induced index $e$ inside the
anchor subgroup, then $nd=4e$; every box has reciprocal capacity one; and
induced indices sharing a box have gcd greater than one. A second exact search
rejects all five profiles. The Lean~4 endpoint has no external theorem
parameter: the finite-quotient reduction, the two partition identities, the
four source-shaped Margolis--Schnabel obstruction arguments, both
group-to-search bridges, and both finite searches are proved in its local
dependency chain. The finite computations are split into checked-in
\texttt{decide +kernel} certificates, and the final axiom report contains only
\texttt{propext}, \texttt{Classical.choice}, and \texttt{Quot.sound}. Two
separately implemented exact-arithmetic Python programs reproduce the
supplementary arithmetic and fiber boundaries; they share the same
mathematical obstruction specification but import no code from one another.
At length eighteen,
39 necessary-condition fiber profiles survive; they are not coset partitions
or counterexamples. To our knowledge, the previously published unrestricted
cell-count result covered at most seven cells.
\end{abstract}

\section{Introduction and statement}

Let $G$ be a group. A finite family of left cosets
\[
  \mathcal P=\{g_iH_i:1\leq i\leq r\}
  \tag{1.1}
\]
of finite-index subgroups is a \emph{coset partition} if every element of $G$
belongs to exactly one displayed coset. Herzog and Sch\"onheim posed the
conjecture in 1974~\cite{HS74}: a nontrivial such partition, meaning $r\geq2$,
cannot have all indices $[G:H_i]$ distinct. See
Margolis--Schnabel~\cite{MS19} and the maintained problem
record~\cite{Erdos274} for background.

The maintained Erd\H{o}s Problems record currently adopts this arbitrary-group
formulation as Problem~274. The historically asked abelian special case is
already known through the subnormal-subgroup result recorded there. Our claim
is a bounded result for the current general-group formulation and is not a new
result for the abelian case.

Our candidate partial theorem is the following.

\begin{theorem}\label{thm:main}
Let $G$ be an arbitrary group and let \emph{(1.1)} be a coset partition with
$2\leq r\leq17$. Then there are distinct $i,j$ such that
\[
  [G:H_i]=[G:H_j].
\]
Equivalently, every counterexample to the Herzog--Sch\"onheim conjecture
requires at least eighteen cells.
\end{theorem}

Akman and Sissokho proved the corresponding arbitrary-group result for at
most seven distinct subgroups~\cite[Theorem~6]{AS25}. In a counterexample the
indices are distinct, so the same subgroup cannot occur twice. Thus their
number of distinct subgroups is exactly the number of cells, and the two
formulations are directly comparable.

The enumeration used here imposes the reciprocal-sum and pairwise-gcd
conditions together with four non-harmonic patterns from
Margolis--Schnabel Propositions 4.2, 4.3, 4.5, and 4.7~\cite{MS19}, and
examines the full remaining denominator interval. The five profiles left at
seventeen are then subjected to the index-four fiber obstruction developed
below.

The theorem is computer assisted, but its trust layers are separated:
\begin{enumerate}[label=(\roman*)]
\item cited articles provide provenance and statement comparison;
\item Lean proves the finite quotient, partition identities, source-shaped
      obstruction theorems, detector transfers, and group-to-search bridges;
\item the arithmetic and fiber computations are divided into finite
      \texttt{decide +kernel} certificates whose coverage is checked by Lean;
\item deterministic generators reproduce every certificate module byte for
      byte, but are not premises of the theorem;
\item implementation-diverse Python programs separately reproduce the
      supplementary length-seventeen and length-eighteen diagnostics.
\end{enumerate}
Section~\ref{sec:trust} states precisely what is and is not formalized.

\section{Finite quotient and index-two descent}\label{sec:group}

We first record two reductions that are easy to hide when moving between the
literature and the executable search.

\subsection{Reduction to a finite group}

\begin{lemma}[Finite quotient]\label{lem:finite-quotient}
Let \emph{(1.1)} be a finite coset partition of an arbitrary group, with every
$H_i$ of finite index. There is a finite quotient $G/N$ carrying a coset
partition with the same number of cells and the same list of indices.
\end{lemma}

\begin{proof}
For each $i$, the normal core
\[
  \operatorname{core}_G(H_i)=\bigcap_{g\in G}gH_ig^{-1}
\]
is the kernel of the action of $G$ on the finite set $G/H_i$, hence is normal
and has finite index. Put
\[
  N=\bigcap_{i=1}^r\operatorname{core}_G(H_i).
\]
Then $N$ is normal and finite index, and $N\leq H_i$ for every $i$. Each cell
$g_iH_i$ is a union of $N$-cosets, so the cells descend to a partition
\[
  \{(g_iN)(H_i/N):1\leq i\leq r\}
\]
of the finite group $G/N$. More explicitly, for the quotient map $\pi$,
\[
  \pi^{-1}\bigl((g_iN)(H_i/N)\bigr)=g_iH_i.
\]
Thus coverage, disjointness, and the number of cells are all preserved.
Finally,
\[
  [G/N:H_i/N]=[G:H_i].
\]
\end{proof}

This is the standard finite reduction in the subject; see also
Korec--Zn\'am~\cite{KZ77}. It permits finite-group harmonic propositions to
be used without changing the profile. The construction, exact cell preimages,
partition descent, and complete index-profile preservation are formalized in
\leanname{FiniteQuotientBridge}.

\subsection{Descent past an index-two cell}

\begin{lemma}[Index-two descent]\label{lem:index-two}
Suppose $\mathcal P$ is a coset partition with $r\geq2$ and all indices are
distinct, and suppose one cell is a coset of an index-two subgroup $K$.
Then $r>2$, and $K$ has a distinct-index coset partition with $r-1$ cells.
\end{lemma}

\begin{proof}
Let the index-two subgroup be $K$. Translate the entire partition so that its
cell is $K$ itself. An index-two subgroup is normal, and the complement of
$K$ is one coset $aK$. Every other cell $g_iH_i$ is disjoint from $K$, hence
is contained in $aK$.

For $h\in H_i$, both $g_i$ and $g_ih$ lie in $aK$. It follows that
$h=g_i^{-1}(g_ih)$ lies in $K$, so $H_i\leq K$. Left translation by $a^{-1}$
therefore sends all cells other than $K$ to a partition of $K$ into cosets of
the same subgroups $H_i$. Index multiplicativity gives
\[
  [G:H_i]=[G:K][K:H_i]=2[K:H_i].
  \tag{2.1}
\]
Thus dividing all remaining indices by two preserves their distinctness.

If $r=2$, the translated family has one cell, so its only subgroup equals
$K$ and has index one in $K$. Equation~(2.1) makes the other original index
equal to two, the same as the removed cell, contrary to distinctness. Hence
$r>2$, and the translated family is a nontrivial distinct-index partition of
$K$ with $r-1$ cells.
\end{proof}

Now suppose a counterexample to Theorem~\ref{thm:main} exists, and choose one
with the fewest cells among all counterexamples having at most seventeen cells.
Lemma~\ref{lem:index-two} shows that none of its indices is two. No minimality
in the group order is needed.

\section{The arithmetic index profile}\label{sec:profile}

Write the distinct indices of the chosen minimal-cell counterexample in
increasing order:
\[
  3\leq a_1<a_2<\cdots<a_r.
  \tag{3.1}
\]
The lower bound follows from the index-two descent and the fact that no cell
of a nontrivial partition can have index one.

Two standard necessary conditions are recorded in
Margolis--Schnabel Lemma 2.3(b,c)~\cite{MS19}:
\begin{align}
  \sum_{i=1}^r\frac1{a_i}&=1, \tag{3.2}\\
  \gcd(a_i,a_j)&>1 \qquad(i\ne j). \tag{3.3}
\end{align}
For a finite quotient, (3.2) is also the identity obtained by dividing the
cardinalities of the partition cells by $|G|$. Condition (3.3) is the familiar
fact that cosets of subgroups with coprime finite indices cannot be disjoint.

We call a list satisfying (3.1)--(3.3), together with the four exclusions in
Section~\ref{sec:harmonic}, an \emph{arithmetic candidate}. Every bounded
counterexample produces one, but the converse is neither used nor expected.

\section{Published harmonic obstructions}\label{sec:harmonic}

A tuple $(d_1,\ldots,d_s)$ is $G$-harmonic if there are subgroups of $G$ with
those indices and representatives whose chosen cosets are pairwise disjoint.
Every selected subfamily of a coset partition is therefore $G$-harmonic.

The following four forbidden shapes are consequences of
Margolis--Schnabel Propositions 4.2, 4.3, 4.5, and 4.7~\cite{MS19}. In each
row, all displayed residual variables are pairwise coprime.

\begin{center}
\begin{tabular}{@{}lll@{}}
\toprule
Source & Forbidden index subtuple & Extra condition \\
\midrule
Prop. 4.2 & $(2r_1,2r_2,2r_3)$ & none \\
Prop. 4.3 & $(3r_1,3r_2,3r_3,3r_4)$ & none \\
Prop. 4.5 & $(2r_1,4r_2,4r_3,4r_4)$ & $r_1$ odd \\
Prop. 4.7 & $(3,3r_2,6r_3,6r_4,6r_5)$ & $r_2$ odd \\
\bottomrule
\end{tabular}
\end{center}

The executable detectors use congruence classes and divide by the common
factor. For example, the Proposition 4.5 detector chooses one index congruent
to $2$ modulo $4$, divides it by $2$, chooses three indices divisible by $4$,
divides them by $4$, and tests pairwise coprimality of the four residuals.
The Proposition 4.7 detector requires the distinguished index $3$, a second
index congruent to $3$ modulo $6$ but unequal to $3$, and three indices
divisible by $6$.

\begin{lemma}[Detector-to-subfamily transfer]\label{lem:detector-transfer}
If the sorted index profile of a distinct-index coset partition triggers one
of the four detectors, then the partition contains a harmonic subfamily with
the corresponding forbidden index tuple.
\end{lemma}

\begin{proof}
Every selected index value comes from a unique cell because the index map is
injective. Selecting the corresponding cells preserves pairwise disjointness.
For Propositions 4.2 and 4.3, the common-factor division gives the residuals
in the source statements. For Proposition 4.5, the index congruent to $2$
modulo $4$ cannot equal any of the three indices divisible by $4$, and its
quotient by $2$ is odd. For Proposition 4.7, the distinguished index $3$, the
nontrivial index congruent to $3$ modulo $6$, and the three multiples of $6$
are all distinct; dividing the second index by $3$ gives an odd residual.
Strict increase handles distinctness within each pool. The source notion is
unchanged by reordering the selected cells. The detector's pairwise
coprimality test is exactly the remaining source hypothesis.
\end{proof}

It follows that the profile of a counterexample triggers none of the four
detectors. Lemma~\ref{lem:finite-quotient} supplies the finite setting if the
original ambient group was infinite.

The local formal proofs use the source-shaped statements but do not copy the
published proofs line by line. Two presentation points are made explicit in
the Lean modules. First, the final inference printed in Corollary~3.11 is
replaced by a direct parity and finite-cardinality argument in the
Proposition~4.3 proof. Second, the penultimate paragraph of the printed
Lemma~4.6 repeats the vector $(1,2,2,3,1,1)$; consistency with the listed
remaining cases requires that paragraph-start occurrence to be
$(1,2,2,3,3,1)$, distinct from the final survivor
$(1,2,2,3,1,1)$. The Proposition~4.7 module records and checks this case split
explicitly.

\section{Complete exact enumeration}\label{sec:search}

We enumerate strictly increasing denominator lists with exact rational
arithmetic. At a recursive node, suppose the positive residual reciprocal sum
is $q$, the previous denominator is $p$, and $k$ denominators remain. If $d$
is the next denominator, then
\[
  \max\{p+1,\lceil 1/q\rceil\}
  \leq d\leq
  \lfloor k/q\rfloor.
  \tag{5.1}
\]

At the root we also impose $d\geq3$.

The lower bound follows from strict increase and $1/d\leq q$. For the upper
bound, every remaining denominator is at least $d$, so their reciprocal sum
is at most $k/d$. Hence $q\leq k/d$.

For every integer $d$ in (5.1), the recursion first performs two monotone
prefix tests:
\begin{enumerate}[label=(\arabic*)]
\item $\gcd(d,a_i)>1$ for every earlier denominator;
\item the extended prefix contains none of the four published obstructions;
\end{enumerate}
It then forms the updated residual $q'=q-1/d$, which is nonnegative by the
lower bound in (5.1), and stores it as a reduced numerator--denominator pair.
A zero residual is accepted only when no denominators remain; if terms remain,
that child state is rejected immediately. Coprime pairs and forbidden
subtuples persist under extension, so both prefix rejection tests are
monotone.

\begin{proposition}[Search completeness]\label{prop:complete}
Every arithmetic candidate of length $r$ is emitted by the exact recursive
search at length $r$.
\end{proposition}

\begin{proof}
Induct on the number of denominators remaining. The next denominator of a
candidate lies in the complete interval (5.1). Conditions (3.3) and the
absence of a published obstruction make it pass both monotone guards. Exact
subtraction sends the candidate tail to the reduced state representing its
reciprocal sum, and strict increase supplies the next previous-denominator
invariant. The induction hypothesis then emits the tail. At zero remaining
terms, equation (3.2) makes the residual exactly zero, so the completed list
is returned.
\end{proof}

This induction is formalized in Lean as \leanname{searchComplete}. The proof
also establishes the natural-number ceiling and floor bounds used by the
implementation and the equality between the rational reciprocal sum and the
reduced integer state.

The length-seventeen equality is not a monolithic evaluator call. The complete
1,052-state recursion tree is divided by depth across
\leanname{ArithmeticSearchCertificateD0} through
\leanname{ArithmeticSearchCertificateD17}. A terminal or local admissible-
denominator check is discharged by \texttt{decide +kernel}; an internal node
unfolds one recursion step and rewrites by every imported child theorem. The
root theorem is \leanname{checked_seventeen_survivors}.

\section{The arithmetic boundary and index-four fiber obstruction}
\label{sec:boundary}

\subsection{The five arithmetic profiles}

The primary arithmetic program uses Python's exact
\texttt{fractions.Fraction} arithmetic and recursive traversal. The
separately implemented verifier stores a reduced pair of natural numbers and
uses an explicit stack; it has its own code for obstruction logic, constants,
hashing, and checked ranges. The two programs import no code from one another
but share the same mathematical obstruction specification. They agree on every
recorded count and endpoint hash.

\begin{center}
\begin{tabular}{@{}lccc@{}}
\toprule
Obstruction profile & Zero through & First survivor length & Survivors \\
\midrule
Propositions 4.2 and 4.3 & $14$ & $15$ & $28$ \\
Propositions 4.2, 4.3, 4.5, 4.7 & $16$ & $17$ & $5$ \\
\bottomrule
\end{tabular}
\end{center}

For the full profile, the primary recursion visits 117, 161, 244, 406, and
1052 nodes at lengths 13, 14, 15, 16, and 17. These implementation-specific
counts are retained only as reproducibility checks.

All five length-seventeen survivors share the prefix
\[
  (4,6,8,12,16,18,24,30,32,36,40,42,48)
\]
and have the following four-term tails:
\[
\begin{array}{@{}cc@{}}
(54,56,60,4320), & (56,60,64,320),\\
(56,60,80,160), & (56,60,96,120),\\
\multicolumn{2}{c}{(56,72,90,96).}
\end{array}
\tag{6.1}
\]
The SHA-256 of their canonical sorted JSON encoding is
\begin{center}
\small\nolinkurl{0244dcaa451200177fe26b8ea2b8ab53f3737817004129afbd9969e1bead2d57}.
\end{center}
In particular, every profile contains the index four. This is a checked fact
about (6.1), not a general assertion about coset partitions.

\subsection{The three-box assignment}

Suppose a distinct-index coset partition has one of the profiles in (6.1).
Choose its unique index-four cell and translate the partition so that this
cell is the subgroup $H$ itself. No normality of $H$ is assumed. The set of
left $H$-cosets has four elements, and the complement of the anchor therefore
consists of three left-coset \emph{boxes}, which we label $1,2,3$.

Consider another cell $gK$ and write $n=[G:K]$. Let $S(gK)$ be the nonempty
set of boxes met by $gK$, put $d=|S(gK)|$, and set
\[
  e=[H:H\cap K].
\]
The set of left cosets of $H\cap K$ in $K$ injects into the set of left
$H$-cosets, with image exactly the boxes met by $gK$. Hence it is in
bijection with that support, so $d=[K:H\cap K]$. Computing
$[G:H\cap K]$ in the two subgroup chains gives
\[
  n d = 4e. \tag{6.2}
\]

If two distinct cells meet the same box, translating that box back to $H$
produces disjoint left cosets of the subgroups $H\cap K$ and $H\cap L$ in
$H$. Cosets of finite-index subgroups with coprime indices always intersect:
indeed, the natural map from one subgroup to the coset space of the other is
surjective when the two indices are coprime. Consequently
\[
  \gcd(e_K,e_L)>1 \tag{6.3}
\]
whenever the two rows share a box.

We also need the exact reciprocal capacity of each box. We use the elementary
coset-cover inequality: if finitely many left cosets of finite-index
subgroups cover a group $L$, then the sum of their reciprocal indices in $L$
is at least one. To see this, pass to the finite quotient by the intersection
of their normal cores and compare cardinalities.

For a fixed box $b$, the fragments of all cells meeting $b$ cover it.
Translating the box to $H$ and applying the cover inequality gives
\[
  C_b:=\sum_{gK:\,b\in S(gK)}\frac1{e_K}\geq1. \tag{6.4}
\]
On the other hand, (6.2) and the ambient reciprocal identity give
\[
  \sum_{b=1}^3 C_b
   =\sum_{gK\ne H}\frac{d_K}{e_K}
   =4\sum_{gK\ne H}\frac1{[G:K]}
   =4\left(1-\frac14\right)=3. \tag{6.5}
\]
There are three boxes and each capacity is at least one, so
\[
  C_1=C_2=C_3=1. \tag{6.6}
\]

Equations (6.2), (6.3), and (6.6) define the finite
\emph{index-four fiber assignment} forced by the partition: every non-anchor
index $n$ chooses a nonempty support $S\subseteq\{1,2,3\}$, with
$e=n|S|/4$ integral; every box has reciprocal $e$-sum one; and all induced
indices sharing a box have gcd greater than one.

\subsection{Kernel depth-first search and split certificates}

We now specify the finite model independently of its implementation. Write a
profile as
\[
  x=(4,n_1,\ldots,n_s),\qquad B=\{0,1,2\}.
\]
For one row $n$, the complete choice set is
\[
  \mathcal C(n)=
  \{(S,e):\varnothing\ne S\subseteq B,\ n|S|=4e\}. \tag{6.7}
\]
Thus the seven nonzero three-bit masks enumerate all possible supports, and a
mask is retained precisely when $n|S|$ is divisible by four. Let
\[
  L=\operatorname{lcm}\{e:(S,e)\in\mathcal C(n_i)
      \text{ for some }i\}. \tag{6.8}
\]
All the profiles under consideration have positive rows, so every displayed
$e$ is positive and divides $L$.

After $t$ rows, a labelled search state is
\[
  Q=((f_0,M_0),(f_1,M_1),(f_2,M_2)). \tag{6.9}
\]
Here $f_b$ is the integer-scaled fill already placed in box $b$, and $M_b$ is
the list, with multiplicity, of induced indices already placed there:
\[
  f_b=\sum_{\substack{i\leq t\\b\in S_i}}\frac{L}{e_i}.
  \tag{6.10}
\]
The initial state is $((0,[]),(0,[]),(0,[]))$. For a choice $(S,e)$ in the
next row, every $b\in S$ is updated by
\[
  (f_b,M_b)\longmapsto
  \left(f_b+\frac Le,e::M_b\right), \tag{6.11}
\]
and the other columns are unchanged. A branch is kept exactly when $e>0$,
$e\mid L$, every new fill is at most $L$, and
\[
  \gcd(e,u)>1\quad\text{for every }u\in M_b\text{ and every }b\in S.
  \tag{6.12}
\]
No quotient by a permutation of the boxes and no state deduplication is used
in the theorem-facing checker. It is a direct depth-first search. After all
rows are processed, a branch accepts exactly when $f_0=f_1=f_2=L$.

The executable procedure can therefore be summarized as follows.

\begin{quote}
\small
\begin{verbatim}
KERNEL-FIBER-DFS(4, n[1], ..., n[s]):
    choices[n] := all nonempty masks S with n*|S| = 4*e
    L := lcm of every induced e in every choices[n]
    DFS(rows, Q):
        if rows is empty:
            return every column fill equals L
        n := first row
        for (S, e) in choices[n]:
            reject if e = 0 or e does not divide L
            reject if a selected fill would exceed L
            reject if gcd(e, u) = 1 for an old u in that box
            if DFS(remaining rows, updated Q): return true
        return false

    return DFS(all rows, ((0, []), (0, []), (0, [])))
\end{verbatim}
\end{quote}

Each rejection is safe for a declarative assignment forced by a partition.
The divisibility guard follows from (6.8). A true prefix cannot overfill,
because all increments are nonnegative and its final fill is exactly $L$ by
(6.6). The gcd guard is exactly (6.3). Finally, after all rows, an underfilled
state cannot satisfy the exact capacity identity. These are the only
rejections used by the theorem-facing Lean transition system.

The search terminates because there are finitely many rows, each row has at
most seven choices, and the row iterator is structural recursion on the finite
tail. Its completeness is the direction needed in the proof.

\begin{proposition}[Depth-first-search completeness]
\label{prop:fiber-search-complete}
Every index-four fiber assignment for $x$ determines a path from the initial
state to an accepting branch of the depth-first search.
\end{proposition}

\begin{proof}
Fix an assignment and process its rows in profile-tail order. Its support in
the next row has a three-bit encoding in $\mathcal C(n)$ by (6.2). The induced
index divides $L$. For every box, the fill of the processed prefix is at most
the final fill $L$, and every old induced index in the same box satisfies
(6.3). Hence the assignment's next transition passes every guard. Induction
on the processed prefix therefore supplies one recursive branch. At the
terminal depth, (6.6) makes all three fills equal to $L$, so the branch
accepts.
\end{proof}

This proposition asserts only
\[
  \text{fiber assignment forced by a partition}
  \Longrightarrow \text{accepted DFS branch}. \tag{6.13}
\]
The reverse construction of a group partition from an accepted branch is not
claimed and is not needed. Therefore an empty executable search rules out all
group-theoretic assignments, whereas a surviving state is only a necessary-
condition object.

The correspondence with the Lean development is explicit.

\begin{center}
\small
\begin{tabular}{@{}p{0.39\linewidth}p{0.53\linewidth}@{}}
\toprule
Mathematical object or step & Lean identifier \\
\midrule
Declarative assignment & \leanname{IndexFourFiberAssignment} \\
Ordered mask witness & \leanname{OrderedFiberMaskAssignment4} \\
Group assignment to ordered witness &
  \leanname{IndexFourFiberAssignment.toOrderedMaskAssignment4} \\
Complete row choices & \leanname{kernelFiberChoices4} \\
Column and three-column update & \leanname{kernelUpdateFiberColumn},
  \leanname{kernelUpdateFiberColumns} \\
Depth-first search and common scale & \leanname{kernelFiberDFS4},
  \leanname{kernelFiberScale4} \\
Assignment-to-DFS completeness &
  \leanname{kernelFiberDFS4_complete_of_orderedAssignment} \\
Five negative profile endpoints & \leanname{kernelFiberProfile1_unsat}
  through \leanname{kernelFiberProfile5_unsat} \\
Combined obstruction &
  \leanname{seventeenIndexFourAssignmentObstruction_kernel} \\
\bottomrule
\end{tabular}
\end{center}

The monolithic closed computation is too large for practical kernel reduction,
so its finite recursion tree is split. A generated leaf contains at most 6,000
visited DFS states and is discharged by \texttt{decide +kernel}. An internal
node unfolds one row, simplifies all concrete guards, and rewrites by every
imported child theorem. Table~\ref{tab:kernel-certs} records the fixed split.

\begin{table}[ht]
\centering
\small
\begin{tabular}{@{}c p{0.31\linewidth} rrr@{}}
\toprule
No. & Four-term tail in (6.1) & $L$ & Leaves & Result \\
\midrule
1 & $(54,56,60,4320)$ & $45{,}360$ & $30$ & false \\
2 & $(56,60,64,320)$  & $30{,}240$ & $30$ & false \\
3 & $(56,60,80,160)$  & $15{,}120$ & $30$ & false \\
4 & $(56,60,96,120)$  & $15{,}120$ & $42$ & false \\
5 & $(56,72,90,96)$   & $15{,}120$ & $42$ & false \\
\bottomrule
\end{tabular}
\caption{Split kernel DFS certificates for the five profiles.}
\label{tab:kernel-certs}
\end{table}

There are 251 generated fiber modules in total: one fixed profile-data module,
174 leaves, 71 internal nodes, and five profile endpoints. The deterministic
generator reconstructs the five trees and reproduces every module byte for
byte. This generator chooses the split, but is not a proof premise: each leaf,
each parent-to-child connection, and each profile endpoint is elaborated by
Lean.

The supplementary primary Python program is deliberately separate from the
Lean proof route. It reorders rows, quotients by box permutations, stores seen
indices as sets, and applies a safe remaining-mass rejection. The separately
implemented verifier keeps boxes labelled and uses bitsets and a different
traversal. Neither implementation is imported by the final Lean theorem.

The optimized primary program and the separately implemented verifier agree
on the following decision boundary. This implementation diversity is a
reproducibility check, not an independent mathematical proof of the model.

\begin{center}
\begin{tabular}{@{}rrrr@{}}
\toprule
Length & Arithmetic profiles & Fiber survivors & Survivor SHA-256 prefix \\
\midrule
$17$ & $5$ & $0$ & \texttt{4f53cda18c2baa0c} \\
$18$ & $470$ & $39$ & \texttt{589513a618402b33} \\
\bottomrule
\end{tabular}
\end{center}

For the empty list at seventeen, the displayed hash is the SHA-256 of the
canonical JSON encoding \texttt{[]}. The complete hashes and all thirty-nine
length-eighteen lists are stored in the JSON artifact.

Proposition~\ref{prop:fiber-search-complete} is formalized as
\leanname{kernelFiberDFS4_complete_of_orderedAssignment}, after
\leanname{IndexFourFiberAssignment.toOrderedMaskAssignment4}. Thus the
negative DFS certificates cannot miss an assignment forced by a partition.

\begin{proposition}[Fiber rejection]\label{prop:fiber-rejection}
None of the five profiles in \emph{(6.1)} admits an index-four fiber
assignment.
\end{proposition}

\begin{proof}
The Lean completeness theorem maps every such assignment to a branch on which
\leanname{kernelFiberDFS4} returns true. The five split profile certificates
prove that the same root searches return false. Their combination is
\leanname{seventeenIndexFourAssignmentObstruction_kernel}; hence no assignment
exists.
\end{proof}

\begin{proof}[Proof of Theorem~\ref{thm:main}]
Assume a counterexample with at most seventeen cells exists and choose one
with the fewest cells. Lemma~\ref{lem:index-two} and minimality first exclude
index two. Apply Lemma~\ref{lem:finite-quotient} and replace the partition by
its finite quotient with the same cell count and complete index profile. Sort
the still-distinct indices. Equations (3.2) and (3.3), together with
Lemma~\ref{lem:detector-transfer} and the four published propositions, make
this list an arithmetic candidate.

If its length is at most sixteen, Proposition~\ref{prop:complete} and the
pure-kernel empty-search computation give a contradiction. The remaining case
has length seventeen. Exact evaluation places its profile among the five in
(6.1). The index-four cell then constructs the three-box assignment above,
contradicting Proposition~\ref{prop:fiber-rejection}.
\end{proof}

The 39 length-eighteen survivors satisfy only the recorded arithmetic and
first-order fiber conditions. They neither realize coset partitions nor show
that an eighteen-cell counterexample exists.

\section{Lean formalization and trust boundary}\label{sec:trust}

The repository uses Lean 4.30.0 and mathlib 4.30.0. The main local endpoints
are:
\begin{description}[style=nextline,leftmargin=8mm]
\item[\leanname{checked_seventeen_survivors}]
the exact five-profile arithmetic boundary from the split D0--D17
certificates;
\item[\leanname{localMargolisSchnabelFacts}]
the locally proved six-field obstruction bundle;
\item[\leanname{kernelFiberDFS4_complete_of_orderedAssignment}]
completeness of the direct labelled depth-first search;
\item[\leanname{seventeenIndexFourAssignmentObstruction_kernel}]
the combined rejection of all five profiles;
\item[\leanname{erdos274AtMostSeventeen}]
the closed at-most-seventeen theorem with no external theorem parameter.
\end{description}

The exact source interface \leanname{MargolisSchnabelFacts} has six fields:
Lemma 2.3(b,c) and detector-level corollaries of Propositions 4.2, 4.3, 4.5,
and 4.7. It remains useful as an internal theorem map, but
\leanname{localMargolisSchnabelFacts} proves every field locally. The final
theorem takes no argument of this type.

The dependency chain also contains the normal-core finite quotient, exact cell
preimages, partition descent, ambient and relative index preservation,
harmonic subfamily selection, and detector-to-cell transfer. Thus no cited
mathematical theorem is left as an external parameter of the bounded endpoint.
This is not a claim that the published article has been transcribed line by
line; the local modules prove the required source-shaped statements directly.

Both theorem-facing finite computations use split \texttt{decide +kernel}
certificates. The arithmetic family has 18 depth modules and the fiber family
has 251 generated modules. The parent theorems prove coverage by unfolding one
transition and rewriting every successful child; a generator is not imported
as an oracle.

For every displayed endpoint, including
\leanname{erdos274AtMostSeventeen}, \texttt{\#print axioms} reports only
\leanname{propext}, \leanname{Classical.choice}, and
\leanname{Quot.sound}. The release checker verifies this allowlist, the final
closed signature with no external theorem parameter, and the absence of
\texttt{sorry}, custom axioms,
\texttt{native\_decide}, \texttt{bv\_decide}, native compiler equality,
\texttt{unsafe}, \texttt{extern}, and \texttt{implemented\_by} in the
first-party theorem tree.

\section{Literature-search scope and priority boundary}
\label{sec:literature-boundary}

Margolis--Schnabel use the four harmonic obstructions in a computation
organized by group order and prove the conjecture for groups of order below
1440~\cite{MS19}. They do not state the complete cell-count enumeration or
index-four fiber obstruction through seventeen.

Akman--Sissokho prove the arbitrary-group result through seven cells and
report that their computation did not finish at eight~\cite{AS25}. Their
later Steiner paper treats structural hypotheses rather than a larger
unrestricted cell-count bound~\cite{ASSteiner,ASSteinerErratum}. Recent work
on simple and symmetric groups also imposes group-class
assumptions~\cite{GM25}.

A bounded public-source search found no earlier arbitrary-group theorem for
eight through seventeen cells. The search covered the maintained problem page,
the primary papers above, public bibliographic indexes, and public citation
graphs, but it was not exhaustive. In particular, no claim is made about all
MathSciNet records, theses, non-English sources, private notes, or folklore.
The defensible wording is therefore:

\begin{quote}
We present an unreviewed computer-assisted proof candidate that every
nontrivial coset partition of an arbitrary group into at most seventeen
cosets has multiplicity. To our knowledge, the previously published
unrestricted bound was seven, due to Akman and Sissokho.
\end{quote}

This qualifies the literature comparison, not the mathematical statement of
Theorem~\ref{thm:main}. The bounded search does not establish priority.

\section{Reproducibility and disclosure}

The standalone repository contains the complete Lean import closure, the two
deterministic kernel-certificate generators, and two pairs of separately
implemented exact-arithmetic Python diagnostics. The paired diagnostic
programs share the same mathematical specification but import no code from
one another. The repository also contains two frozen JSON artifacts, short
source records, this manuscript, and SHA-256 manifests. The command
\begin{center}
\texttt{bash scripts/check\_release.sh}
\end{center}
runs the artifact replays, byte-for-byte certificate checks, a two-job
controlled Lean build, theorem and axiom audit, source scan, and manifest
checks. A release decision should also include a build from a fresh copy
without local Lean artifacts.

LLMs accessed through ChatGPT Pro and OpenAI Codex, primarily GPT-5.6, were
used in proof exploration, formalization, test generation, literature
screening, and editing. Rio Itabe is responsible for the mathematical claim,
source use, code, manuscript, and release decisions. LLM output was not
treated as a proof, bibliographic authority, or independent review.

\begin{thebibliography}{9}
\small

\bibitem{AS25}
F. Akman and P. A. Sissokho,
\emph{Transversal coset partitions of groups},
Beitr. Algebra Geom. \textbf{66} (2025), 417--441.
\href{https://doi.org/10.1007/s13366-024-00748-9}
{doi:10.1007/s13366-024-00748-9}.

\bibitem{ASSteiner}
F. Akman and P. A. Sissokho,
\emph{Steiner coset partitions of groups},
Canad. Math. Bull. \textbf{68} (2025), no.~4, 1359--1373.
\href{https://doi.org/10.4153/S0008439525100787}
{doi:10.4153/S0008439525100787}.

\bibitem{ASSteinerErratum}
F. Akman and P. A. Sissokho,
\emph{Steiner coset partitions of groups -- erratum},
Canad. Math. Bull. \textbf{69} (2026), no.~2, 702.
\href{https://doi.org/10.4153/S0008439525101392}
{doi:10.4153/S0008439525101392}.

\bibitem{Erdos274}
T. F. Bloom,
\emph{Erd\H{o}s Problem 274}, maintained problem record,
accessed 16 August 2026.
\href{https://www.erdosproblems.com/274}{erdosproblems.com/274}.

\bibitem{GM25}
M. Garonzi and L. Margolis,
\emph{The Herzog--Sch\"onheim conjecture for simple and symmetric groups},
arXiv:2509.25118v2 (2025; revised 2026), forthcoming in Algebra \& Number
Theory.
\href{https://arxiv.org/abs/2509.25118}{arXiv:2509.25118}.

\bibitem{HS74}
M. Herzog and J. Sch\"onheim,
\emph{Research problem no.~9},
Canad. Math. Bull. \textbf{17} (1974), 150.

\bibitem{KZ77}
I. Korec and \v{S}. Zn\'am,
\emph{On disjoint covering of groups by their cosets},
Math. Slovaca \textbf{27} (1977), no.~1, 3--7.
\href{https://dml.cz/dmlcz/129020}{dml.cz/dmlcz/129020}.

\bibitem{MS19}
L. Margolis and O. Schnabel,
\emph{The Herzog--Sch\"onheim conjecture for small groups and harmonic
subgroups},
Beitr. Algebra Geom. \textbf{60} (2019), 399--418.
\href{https://doi.org/10.1007/s13366-018-0419-1}
{doi:10.1007/s13366-018-0419-1}.

\end{thebibliography}

\end{document}
